\documentclass[11pt]{article}
\usepackage[a4paper,margin=1in]{geometry}
\usepackage{amsmath,amssymb,amsfonts}
\usepackage{graphicx}
\usepackage{booktabs}
\usepackage{multirow}
\usepackage{siunitx}
\usepackage{enumitem}
% Hyperlinks with stable unique anchors and line-breaking
\usepackage[hidelinks,breaklinks=true,hypertexnames=false]{hyperref}
% Micro-typography to reduce overfull boxes
\usepackage[final]{microtype}
\usepackage[numbers,sort&compress]{natbib}
\usepackage{algorithm}
\usepackage[noend]{algpseudocode}
\usepackage{CJKutf8}
\usepackage{authblk}

\title{\textbf{AI-Philo 1.0: Cross-Disciplinary AI for Automated Synthesis, Evaluation, and Evolution of Quantum Interpretations}}
\author[1]{[Your Name]}
\author[1]{[Co-author Name]}
\affil[1]{[Affiliation], [City], [Country]}
\date{\today}

\begin{document}
\maketitle

% ---------------- Abstract (EN) ----------------
\begin{abstract}
We present \textsc{AI-Philo}, a cross-disciplinary framework that \emph{generates}, \emph{evaluates}, and \emph{evolves} candidate interpretations of quantum mechanics using large language models (LLMs) and human-in-the-loop protocols. 
The method combines contradiction-driven theory synthesis with concept-space exploration over a literature-augmented embedding. 
Evaluation integrates five canonical experiments (Bell/CHSH, single-electron double-slit, delayed-choice quantum eraser, collisional decoherence with fullerenes, Wheeler delayed choice) with role-based expert scoring into a composite metric, and is reinforced by blinded comparative review with Bradley--Terry scaling, pre-registered forecasting with proper scoring rules, MDL-based selection with a minimal ``prepare--evolve--readout'' DSL, and formal consistency checks (no-signalling; relations to Bell/KS/PBR/Gleason). 
Across runs, three AI-generated interpretations---PRQM, CERT, and IEQT---reach a composite score of \textbf{0.907}, approaching Many-Worlds (\textbf{0.913}) and tying Consistent Histories (\textbf{0.907}) under our protocol. 
We release templates, scripts, and anonymized theory cards to support reproducibility.
\end{abstract}

\paragraph{Keywords} Scientific discovery automation; quantum foundations; large language models; theory synthesis; concept spaces; blinded human review; MDL; reproducibility.

% =====================
% 1. Introduction
% =====================
\section{Introduction}
Interpretations of quantum mechanics (QM)---Copenhagen, Many-Worlds, de~Broglie--Bohm, objective collapse, relational QM, consistent histories, and others---feature divergent ontological and epistemic commitments while largely matching standard predictions.
Exhaustively exploring this space is hard: the hypothesis space is vast, the evaluation criteria are multi-faceted, and iterative refinement is expensive.
\textsc{AI-Philo} addresses these challenges with a cross-disciplinary pipeline that couples LLM-based theory generation with auditable human evaluation and reproducible engineering.

\paragraph{Contributions.} This paper makes six contributions:
\begin{sloppypar}
\begin{enumerate}[leftmargin=2em]
  \item \textbf{Contradiction-driven synthesis} that extracts ontological/\allowbreak epistemic/\allowbreak mathematical/\allowbreak empirical tensions across prior interpretations and proposes bridging hypotheses.
  \item \textbf{Concept-space exploration} that identifies low-density gaps in a literature-augmented embedding and generates novel combinations.
  \item \textbf{Composite evaluation} that combines agreement with five canonical experiments~\cite{Aspect1982,Tonomura1989,Kim2000,Hornberger2003,Jacques2007} and role-based expert scoring.
  \item \textbf{Evolutionary refinement} via selection--refinement--mutation across generations.
  \item \textbf{Human-in-the-loop suite} beyond LLM pre-screening: blinded comparative review with Bradley--Terry/\allowbreak Elo scaling; pre-registered forecasting with proper scoring rules (log/\allowbreak Brier); MDL with a minimal ``prepare--evolve--readout'' DSL; and formal checks (no-signalling; Bell/\allowbreak KS/\allowbreak PBR/\allowbreak Gleason).
  \item \textbf{Reproducibility} via repository audit, released templates, scripts, and anonymized theory cards.
\end{enumerate}
\end{sloppypar}

% =====================
% 2. Related Work
% =====================
\section{Related Work}
\paragraph{Quantum interpretations.} Foundational proposals include Everett's relative-state (Many-Worlds)~\cite{Everett1957}, Consistent Histories~\cite{Griffiths1984}, and Relational QM~\cite{Rovelli1996}. Canonical experiments that anchor evaluation include Aspect's Bell tests~\cite{Aspect1982}, Tonomura's single-electron interference~\cite{Tonomura1989}, Kim \emph{et al.}'s delayed-choice quantum eraser~\cite{Kim2000}, Hornberger \emph{et al.}'s collisional decoherence~\cite{Hornberger2003}, and Jacques \emph{et al.}'s Wheeler delayed choice~\cite{Jacques2007}.
\paragraph{AI for scientific discovery.} LLMs and program synthesis have been used to hypothesize mechanisms, unify conceptual spaces, and drive iterative design; our work emphasizes contradiction resolution, gap-finding, and auditable evaluation specific to quantum foundations.

% =====================
% 3. Methods
% =====================
\section{Methods}
\subsection{Theory Generation}
\textbf{Contradiction-driven synthesis.} Given a set of prior interpretations $\mathcal{T}$, we enumerate pairs $(T_i,T_j)$, elicit tensions across ontology/epistemology/mathematics/empirics, and propose bridging hypotheses with internal coherence checks.
\begin{algorithm}[h]
\caption{Contradiction-driven synthesis (sketch)}
\begin{algorithmic}[1]
\Require prior set $\mathcal{T}$
\ForAll{pairs $(T_i,T_j)\in \mathcal{T}\times\mathcal{T}$}
  \State extract tensions; generate hypotheses; validate coherence
\EndFor
\State score $\rightarrow$ select $\rightarrow$ refine $\rightarrow$ next generation
\end{algorithmic}
\end{algorithm}

\textbf{Concept-space exploration.} We embed core concepts into a high-dimensional space (literature-augmented), use clustering/density analysis to identify low-density regions, and generate combinations via interpolation/extrapolation with provenance tracking.

\subsection{Evaluation Protocol}
\textbf{Canonical experiments.} Agreement checks cover Bell/CHSH, single-electron double-slit, delayed-choice quantum eraser, collisional decoherence with fullerenes, and Wheeler delayed choice.\\
\textbf{Role-based scoring.} Three expert roles---physicist, philosopher, mathematician---score mathematical rigor, testability, ontological clarity, and coherence.\\
\textbf{Composite metric.} Unless stated otherwise, we use $0.6\times$ experimental agreement $+$ $0.4\times$ role average.

\subsection{Human-in-the-Loop Robust Evaluation}
We extend beyond LLM pre-screening with six channels: (1) blinded theory cards; (2) Bradley--Terry/Elo scaling; (3) pre-registered forecasting with proper scoring rules (log/Brier); (4) MDL with a minimal ``prepare--evolve--readout'' DSL; (5) formal checks (no-signalling; Bell/KS/PBR/Gleason); (6) MCDM with robustness analysis.

\subsection{Evolutionary Optimization}
A selection--refinement--mutation controller improves candidates over generations while preserving diversity and avoiding premature convergence.

\subsection{Mathematical Relationship to Standard QM}
We summarize alignment and deviations using a compliance map (Table~\ref{tab:compliance}); deviations specify domain, conserved quantities, Born-rule status (postulated vs.\ derived), and no-signalling guarantees.
\begin{table}[t]
\centering
\small
\caption{Compliance map between standard QM and AI-generated interpretations (to be filled with finalized specifics).}
\label{tab:compliance}
\begin{tabular}{lccc}
\toprule
Module & Standard QM & This Work & Relation \\
\midrule
State space & Complex Hilbert space & \textit{[edit]} & Equal/Extended \\
Composition & Tensor product & \textit{[edit]} & Equal/Modified \\
Dynamics & Schr\"odinger / LvN & \textit{[edit]} & Equal/Nonlinear/Stochastic \\
Measurement & POVM / instruments & \textit{[edit]} & Equal/Derived/Replaced \\
Born rule & Postulate & \textit{[edit]} & Kept/Derived/Modified \\
Symmetry & Wigner/Lorentz & \textit{[edit]} & Kept/Weakened/Extended \\
No-signalling & Yes & \textit{[edit]} & Proof/Counterexample \\
Classical limit & $\hbar\!\to\!0$ & \textit{[edit]} & Recovered/Conditions \\
\bottomrule
\end{tabular}
\end{table}

% =====================
% 4. System Overview
% =====================
\section{System Overview}
Figure~\ref{fig:arch} summarizes data flow: prior interpretations $\rightarrow$ contradiction extraction \& concept-space exploration $\rightarrow$ candidate theories $\rightarrow$ evaluation (canonical experiments + role scores + human-in-the-loop channels) $\rightarrow$ evolutionary refinement.
\begin{figure}[h]
\centering
\fbox{\parbox{0.92\linewidth}{\small \textbf{Pipeline.} Prior theories $\rightarrow$ synthesis (contradictions \& concept space) $\rightarrow$ candidates $\rightarrow$ evaluation (experiments, roles, human-in-the-loop) $\rightarrow$ evolutionary refinement.}}
\caption{System architecture (schematic).}
\label{fig:arch}
\end{figure}

% =====================
% 5. Experiments
% =====================
\section{Experiments}
\subsection{Setup}
We evaluate 13 prior interpretations and two generation routes (contradiction-driven, concept-space). Unless otherwise noted, three generations are run with 2--3 survivors per generation.
\subsection{Baselines and Metrics}
Baselines include Many-Worlds, Consistent Histories, Copenhagen, and others. Metrics include canonical experiment agreement, role-based scores, the composite metric, MDL, and forecasting scores.
\subsection{Results}
Across runs, AI-generated interpretations reach a best composite score of \textbf{0.907}; Many-Worlds reaches \textbf{0.913}, Consistent Histories \textbf{0.907}, Copenhagen 0.900. 
Three AI theories---PRQM, CERT, IEQT---tie at 0.907 under our composite metric.
\begin{table}[h]
\centering
\small
\caption{Summary of generation scale and headline scores (update with latest runs if needed).}
\label{tab:summary}
\begin{tabular}{lcccc}
\toprule
Method & \#Theories & Avg.\ Score & Best Score & Success \\
\midrule
Contradiction-driven & $127{+}$ & 0.80 & \textbf{0.907} & 85\% \\
Concept-space        & $89{+}$  & 0.77 & 0.84 & 80\% \\
Overall              & $200{+}$ & 0.79 & \textbf{0.907} & 82\% \\
\bottomrule
\end{tabular}
\end{table}

\subsection{Case Studies}
PRQM, CERT, and IEQT differ in ontology and measurement treatment while preserving no-signalling and reproducing canonical outcomes under our evaluation protocol.

% =====================
% 6. Discussion
% =====================
\section{Discussion}
\paragraph{Strengths.} Systematic exploration via contradictions and concept gaps; auditable human evaluation; MDL and forecasting encourage parsimony and calibration; reproducible pipeline with released assets.
\paragraph{Limitations.} LLMs may propose mathematically descriptive but under-derivable formalisms; role prompts risk bias; forecasting timelines can be long; formalization and QFT extensions require additional work.
\paragraph{Opportunities.} Integrate symbolic derivations and theorem provers; expand to new experimental regimes (macroscopic interference, gravity-related decoherence); apply the workflow to other theory-rich domains.

% =====================
% 7. Conclusion
% =====================
\section{Conclusion}
Contradiction-driven synthesis plus concept-space exploration, reinforced by human-in-the-loop evaluation and reproducible engineering, can traverse the interpretation landscape and repeatedly surface high-quality candidates approaching leading prior theories.

% ---------------- References ----------------
\bibliographystyle{plainnat}
\bibliography{ai-philo-refs}
\end{document}
